The current Linux backing remains \texttt{rw-s}; read-only protection is not enforced after population. Each lane populates the shared arena during startup, and the exact-size gate does not prove model-content identity. Sequential startup avoids concurrent initial population, but a live lane restart can rewrite pages that other lanes map. The deployment therefore assumes identical model contents across lanes, and shared-memory corruption or misconfiguration is a cross-lane failure mode. We make no fault-isolation claim for the weight arena. Finally, the workload study is observational: speculative acceptance, routing locality, and context length are not independently controlled, and no adaptive swaps occur in the retained workload windows.

\section{Reproducibility and Artifact}
The companion repository keeps the runtime implementation and machine-specific evidence deliberately separate. The Strata fork holds generic serving code and contracts, while the recipe repository holds reference-host methodology, raw repetition-level CSVs, current and historical promotion records, and reporting rules~\cite{artifact}. The final evidence pin records exact fork/upstream commits, binary SHA-256, driver/CUDA versions, topology, prompt hashes, and anomaly notes. Historical 0.1.22/0.1.24 measurements remain available but are not treated as strict software-version A/Bs when prompts or measurement contracts differ.

